
en
vertu de 4.5, l'implication pour @@M0@@ en r\'esulte. Supposons donc dor\'enavant @@M1@@ projectif, muni d'un faisceau ample @@M2@@.

Soient @@M3@@ les composantes irr\'eductibles de @@M4@@, munies de la structure r\'eduite induite. Appliquant 4.5 et EGA III 2.6.2 au morphisme fini et surjectif @@M5@@, on ram\`ene la d\'emonstration que (iii) implique (i) et (iv) au cas o\`u @@M6@@ est int\`egre.

Si @@M7@@ est irr\'eductible, il est \'evident que (ii) implique (v), que (iii) implique (vi), et compte tenu de 2.1, que (v) implique (vi), et que (vi) pour @@M8@@ est \'equivalent \`a (vi) pour @@M9@@ (X 4.4.1).

Il reste donc \`a voir que (vi) implique (i) et (iv), et que (iv) implique (i), en supposant @@M10@@ int\`egre. Mais alors, (vi) (resp. (iv)) implique que la famille @@M11@@ des classes des @@M12@@ (@@M13@@) satisfait \`a 3.13 (iii) (resp. (iv)), donc est limit\'ee. Il existe donc un faisceau inversible @@M14@@ sur un @@M15@@, avec @@M16@@ de type fini sur @@M17@@, qui limite @@M18@@. Soient @@M19@@ deux entiers diff\'erents tels que @@M20@@ et @@M21@@ correspondent \`a la m\^eme composante connexe de @@M22@@, alors @@M23@@ est donc alg\'ebriquement \'equivalent \`a @@M24@@\,; d'o\`u (vi) (resp. (iv)) implique (i).

Enfin, on a de m\^eme que (vi) implique qu'il existe un entier @@M25@@ tel que pour tout entier @@M26@@, @@M27@@ soit ample. Alors pour tout @@M28@@, @@M29@@ est ample\,; d'o\`u que (vi) implique (iv), ce qui ach\`eve la d\'emonstration.

\textbf{Th\'eor\`eme 4.7.} \emph{Soient @@M30@@ un sch\'ema noeth\'erien, et @@M31@@ un @@M32@@-sch\'ema propre. Alors}~:
\begin{enumerate}[label=(\roman*)]
\item \emph{Le morphisme d'inclusion @@M33@@ est repr\'esentable par une immersion ouverte.}
\item
\emph{Si @@M34@@ est projectif et int\`egre sur @@M35@@, alors @@M36@@ 
